\Needspace{9\baselineskip}\begin{center}On the average number of genera.\end{center}


\Needspace{7\baselineskip}\Article{301}
We return to binary forms, several of whose remarkable properties still need to be discussed. First we will add some observations on the number of genera and classes in the properly primitive order, positive for negative determinants. For brevity we restrict the investigation to this order.

The \textit{number of genera} into which all properly primitive forms, positive for negative determinants, of a given positive or negative determinant \(\pm D\) divide is always \(1,2,4\) or a higher power of \(2\). Its exponent depends on the factors of \(D\) and can be found entirely in advance from the preceding investigations. Since primes are mixed with numbers of varying degrees of compositeness in the natural series of numbers, the number of genera for several
\SourceLeaf{368}
successive determinants \(\pm D,\pm(D+1),\pm(D+2)\), and so on, now increases and now decreases, and this irregular series seems to have no order. Yet if we add the numbers of genera corresponding to many successive determinants
\[\pm D,\quad\pm(D+1)\;\ldots\;\pm(D+m)\]
and divide the sum by the number of determinants, we obtain the \textit{average number of genera}. This can be regarded as applying around the middle determinant \(\pm(D+\frac12m)\), and forms a very regular progression. We assume not only that \(m\) is large enough, but also that \(D\) is much larger, so that the ratio of the extreme determinants \(D,D+m\) does not differ too much from equality. The regularity of this progression is to be understood as follows. If \(D^\prime\) is much larger than \(D\), the average number of genera around the determinant \(\pm D^\prime\) will be noticeably larger than around \(D\). But if \(D,D^\prime\) do not differ too much, the average numbers of genera around \(D\) and \(D^\prime\) will also be nearly equal. Further, the average number of genera around the positive determinant \(+D\) is always found nearly equal to the average around the negative one, and the more closely so the larger \(D\) is. For a small value the former is a little larger than the latter. The following examples, taken from a table classifying binary forms for more than \(4000\) determinants, will illustrate these observations. Among the hundred determinants from \(801\) to \(900\), there are \(7\) to which just one genus corresponds, and \(32,52,8,1\) to which \(2,4,8.\ 16\) genera correspond, respectively. This gives \(359\) genera in all, so the average number \(=3{,}59\). The hundred negative determinants from \(-801\) to \(900\) give \(360\) genera. All the following examples come from negative determinants. In the \(16\)th hundred, from \(-1501\) to \(-1600\), the average number of genera is \(3{,}89\); in the \(25\)th hundred it is \(4{,}03\); in the \(51\)st it is \(4{,}24\); from the six hundred determinants \(-9401\;\ldots -10000\) we obtain \(4{,}59\). These examples show that the average number of genera grows much more slowly than the determinants themselves. But what is the law of this progression? — A rather difficult theoretical investigation, which would take too long to explain here, has shown that the average number of genera around the determinant \(+D\) or \(-D\) is very closely given by
\[\alpha\log D+\beta\]

\SourceLeaf{369}
where \(\alpha,\beta\) are constants, namely
\[\alpha=\frac{4}{\pi\pi}=0{,}4052847346\]
with \(\pi\) denoting half the circumference of a circle of radius \(1\), and
\[\beta=2\alpha g+3\alpha\alpha h-\frac16\alpha\log2=0{,}8830460462\]
where \(g\) is the sum of the series
\[\begin{aligned}&1-\log(1+1)+\frac12-\log(1+\frac12)\\&\quad+\frac13-\log(1+\frac13)+\text{etc.}=0{,}5772156649\end{aligned}\]
(see Euler, \textit{Inst. Calc. Diff.}, p.\ 444), while \(h\) is the sum of the series
\[\frac14\log2+\frac19\log3+\frac1{16}\log4+\text{etc.}\]
which has been found by approximation \(=0{,}9375482543\). This formula shows that the average number of genera increases in an arithmetic progression when the determinants increase in a geometric progression. The values of this formula for \(D=850\frac12\), \(1550\frac12,2450\frac12,5050\frac12,9700\frac12\) are \(3{,}617;3{,}86;4{,}046;4{,}339;4{,}604\), differing little from the averages given above. The larger the middle determinant, and the more determinants used to calculate the average, the less it will differ from the value of the formula. This formula also lets us find very closely the total number of genera corresponding to successive determinants \(\pm D,\pm(D+1)\ldots\pm(D+m)\), however far apart the extremes \(D,D+m\) are, by calculating and adding the corresponding average numbers for each. This sum will be
\[=\alpha(\log D+\log(D+1)+\text{etc.}+\log(D+m))+\beta(m+1)\]
or, with sufficient accuracy,
\[\begin{aligned}={}&\alpha((D+m)\log(D+m)-(D-1)\log(D-1))\\&+(\beta-\alpha)(m+1)\end{aligned}\]
In this way the sum of the numbers of genera for the determinants \(-1\) to \(-100\) is found \(=234{,}4\), while it is actually \(233\); similarly, from \(-1\) to \(-2000\), \(=7116{,}6\), while it is \(7112\); from \(-9001\) to \(-10000\), where it is \(4595\), the formula gives \(4594{,}9\). Such agreement could hardly have been expected.
\SourceLeaf{370}

\Needspace{9\baselineskip}\begin{center}On the average number of classes.\end{center}


\Needspace{7\baselineskip}\Article{302}
As regards the \textit{number of classes}, always understood as properly primitive classes, positive for negative determinants, positive determinants behave quite differently from negative ones. So we will consider the two separately. They agree in this: for a given determinant, each genus contains equally many classes. Thus the total number of classes equals the product of the number of genera and the number of classes in each genus.

First, for negative determinants, the number of classes corresponding to several successive determinants \(-D,-(D+1),-(D+2)\), and so on, forms just as irregular a progression as the number of genera. But the average number of classes, which will need no definition, grows very regularly, as the following examples show. The hundred determinants from \(-500\) to \(-600\) give \(1729\) classes, so the average number \(=17{,}29\). Similarly, in the \(15\)th hundred the average number of classes is \(28{,}26\); from the two hundreds \(24\) and \(25\) we obtain \(36{,}28\); from the three \(61,62\) and \(63\) we get \(58{,}50\); from the five \(91\ldots95\), \(71{,}56\); finally, from the five \(96\ldots100\), \(73{,}54\). These examples show that the average number of classes grows more slowly than the determinants, but much faster than the average number of genera. A little attention shows that it grows quite closely in proportion to the square roots of the middle determinants. Indeed, a theoretical investigation has shown us that the average number of classes around the determinant \(-D\) is approximately expressed by
\[\gamma\surd D-\delta\]
where
\[\gamma=0{,}7467183115=\frac{2\pi}{7e}\]
with \(e\) denoting the sum of the series
\[1+\frac18+\frac1{27}+\frac1{64}+\frac1{125}\;\text{etc.}\]

\[\delta=0{,}2026423673=\frac2{\pi\pi}\]
The averages calculated by this formula differ little from those taken above from the table of classifications. This formula also lets us find very closely the total number of all properly primitive classes, positive for negative determinants, corresponding to
\SourceLeaf{371}
successive determinants \(-D,-(D+1),-(D+2)\ldots-(D+m-1)\), however far apart the extremes may be, by adding the averages corresponding to those determinants according to the formula. This gives
\[=\gamma(\surd D+\surd(D+1)+\text{etc.}+\surd(D+m-1))-\delta m\]
or, very nearly,
\[=\frac23\gamma((D+m-\frac12)^{\frac32}-(D-\frac12)^{\frac32})-\delta m\]
For example, for the hundred determinants \(-1\ldots-100\) this total is calculated from the formula \(=481{,}1\), while it is actually \(477\). The thousand determinants \(-1\ldots-1000\) give \(15533\) classes according to the table; the formula gives \(15551{,}4\). The second thousand gives \(28595\) classes according to the table, and the formula \(28585{,}7\); similarly, the third thousand actually gives \(37092\) classes, and the formula \(37074{,}3\); the tenth thousand gives \(72549\) by the table, and \(72572\) by the formula.
